% ============================================================
%  §3.2 Θερμικό μοντέλο ελαστικού
% ============================================================

\label{subsec:thermal_model}

Το θερμικό μοντέλο των ελαστικών που αναπτύσσεται στην παρούσα εργασία είναι ένα μοντέλο
συγκεντρωμένων ιδιοτήτων (\textit{lumped}), ενός βαθμού ελευθερίας (θερμοκρασία πέλματος)
και ακολουθεί κατά κύριο λόγο το μοντέλου που αναπτύσσουν οι Tremblett και Limebeer \cite{Tremlett2016}.
Το μοντέλο αυτό θεωρεί πως το πέλμα δε συναλλάσσει θερμότητα με τις γειτονικές του μάζες
και αλληλεπιδρά παρά μόνο μέσω της θερμότητας που παράγεται και απάγεται απο το περιβάλλον.
Στην πράξη αυτή η παραδοχή αποτελεί σημαντική απλούστευση της πραγματικότητας: αν αγνοήσουμε
παράπλευρες πηγές ή καταβόθρες θερμότητας, η μάζα του σκελετού δρα ως μία "δεξαμενή" θερμότητας
που σταθεροποιεί τη θερμοκρασία του πέλματος κατά την οδήγηση. Επιπλέον, λόγω της θερμικής
αντίστασης μεταξύ πέλματος και σκελετού του ελαστικού, το πέλμα έχει την ελευθερία να
θερμανθεί παραπάνω από τον σκελετό του ελαστικού στιγμιαία. Επομένως, δε θα ήταν σωστή
η μοντελοποίηση ολόκληρης της μάζας του ελαστικού με μία μεταβλητή κατάστασης καθώς αυτό θα
οδηγούσε σε πολύ υψηλές χρονικές σταθερές -- δηλαδή πολύ αργή απόκριση του πέλματος
σε θερμότητα \cite{West2022}. Το παραπάνω φαινόμενο είναι η κυριότερη μελλοντική
βελτίωση του μοντέλου ως συγκεντρωμένων ιδιοτήτων, καθώς υπάρχουν στη βιβλιογραφία
και κατανεμημένα θερμικά μοντέλα πεπερασμένων στοιχείων, με κύριο σκοπό την επιπλέον μοντελοποίηση
και της κατανομής θερμοκρασίας στο πέλμα \cite{Farroni2014TRT}. \\[0.2cm]

Η επιλογή του απλούστερου μοντέλου είναι σκόπιμη. Η επίλυση ενός συστήματος
δύο Σ.Δ.Ε εισάγει αμελητέο υπολογιστικό κόστος συγκριτικά με ολόκληρο τον υπολογισμό
της δυναμικής του οχήματος. Ωστόσο, η επιλογή πολυπλοκότερου μοντέλου εισάγει μεγάλο
αριθμό παραμέτρων και φυσικών ιδιοτήτων που αφενός συνεπάγονται μεγαλύτερη ευαισθησία
στο μοντέλο, και αφετέρου χρειάζονται ακόμη μεγαλύτερο σύνολο πειραματικών 
δεδομένων για τον χαρακτηρισμό τους. Χρησιμοποιείται το θερμικό μοντέλο
και οι φυσικές παράμετροι των Tremblett και Limebeer \cite{Tremlett2016} οι οποίοι
επιτυγχάνουν αξιοσημείωτη συσχέτιση με πειραματικά δεδομένα, ακόμη και με απλούστερη
μοντελοποίηση.

\subsubsection{Υπολογισμός θερμότητας}

Στο θερμικό μοντέλο λαμβάνονται υπόψη οι εξής κυρίαρχες μορφές θερμότητας.

\begin{enumerate}
    \item $Q_1$ -- Θερμότητα λόγω τριβής της περιοχής ολίσθησης με το οδόστρωμα
    \item $Q_2$ -- Θερμότητα λόγω υστέρησης του ελαστομερούς
    \item $Q_3$ -- Θερμότητα συναγωγής με περιβάλλων αέρα
    \item $Q_4$ -- Θερμότητα αγωγής με οδόστρωμα
\end{enumerate}

\paragraph{Θερμότητα λόγω τριβής.}

Η ισχύς που παράγεται λόγω τριβής της περιοχής ολίσθησης του πέλματος υπολογίζεται ως
άθροισμα των όρων ισχύος του διαμήκους και εγκάρσιου φορτίου.

\begin{equation}
    Q_1 = p_1u_n\bigl(\lvert F_x\kappa\rvert + \lvert F_ytan(\alpha)\rvert\bigr)
    \label{eq:Q1}
\end{equation}

Όπου $\alpha,\kappa$ τα μεγέθη ολίσθησης, $F_x, F_y$ 
η διαμήκης και εγκάρσια δύναμη, $u_n$ η διαμήκης ταχύτητα του ελαστικού, 
και $p_1$ αδιάστατη παράμετρος που διορθώνει για το ποσό θερμότητας που απορροφά
το οδόστρωμα.

\paragraph{Θερμότητα λόγω υστέρησης.}

Η ισχύς που παράγεται λόγω υστέρησης υπολογίζεται ως
άθροισμα των όρων ισχύος που παράγονται λόγω καθενός από τα διαμήκη, εγκάρσια και κατακόρυφα φορτία.

\begin{equation}
    Q_2 = u_n\bigl(p_2\lvert F_x \rvert
    + p_3\lvert F_y \rvert
    + p_4\lvert F_z \rvert
    \bigr)
    \label{eq:Q2}
\end{equation}

Όπου $F_{\{x,y,z\}}$ το διαμήκες, εγκάρσιο και κατακόρυφο φορτίο αντίστοιχα,
και οι παράμετροι $p_{2\dots4}$ εκφράζουν την απόδοση κάθε όρου -- δηλαδή
πόσο η κυκλική παραμόρφωση που παράγει κάθε όρος μεταφράζεται σε θερμότητα.
\paragraph{Θερμότητα συναγωγής.}

Για τη συναγωγή με τον αέρα περιβάλλοντος θερμοκρασίας $T_{\text{περ.}}$ χρησιμοποιείται ένας
Νευτώνειος νόμος ψύξης.

\begin{equation}
    Q_3 = p_5u^{p_6}\bigl(T_s - T_{\text{περ.}}\bigr)
    \label{eq:Q3}
\end{equation}


Όπου $u$ η ταχύτητα του οχήματος και $p_{5,6}$ παράμετροι που εκφράζουν τη διαφορά στη ροή
αέρα που βλέπει ο κάθε τροχός.

\paragraph{Θερμότητα αγωγής με οδόστρωμα.}

Η θερμότητα που άγεται μεταξύ του ελαστικού στην περιοχή προσκόλλησης και το οδόστρωμα δίνεται από

\begin{equation}
    Q_4 = h_tA_{\text{πέλμα}}\bigl(T_s - T_{οδ.}\bigr)
    \label{eq:Q4}
\end{equation}

Όπου $h_t$ ο συντελεστής συναλλαγής θερμότητας στη διεπαφή ελαστικού -- ασφάλτου και $A_{\text{πέλμα}}$
η επιφάνεια επαφής της περιοχής προσκόλλησης του πέλματος.

Για τον προσδιορισμό των φυσικών μεγεθών μπορούν να ακολουθηθούν διάφορες προσεγγίσεις ή μετρήσεις, ωστόσο,
ακολουθούμε το μοντέλο των Tremblett \& Limebeer \cite{Tremlett2016} και τα προσεγγίζουμε με τις 
\cref{eq:contactArea1,eq:contactArea2,eq:contactArea3}.

Η επιφάνεια επαφής του πέλματος θεωρείται ορθογωνική με πλάτος $C_w$ και μήκος $C_l$. Το πλάτος μπορεί 
να διαφέρει μεταξύ των εμπρόσθιων και οπίσθιων τροχών, οπότε το χωρίζουμε σε $C_{wf},C_{wr}$ για πλάτος
των εμπρός και πίσω ελαστικών αντίστοιχα.
Το μήκος της επιφάνειας επαφής του πέλματος δίνεται από την εμπειρική εκθετική σχέση 
(\ref{eq:contactArea1}).
\begin{equation}
    C_l = \alpha_{cp}F_z^{0.7}
    \label{eq:contactArea1}
\end{equation}

Η εμπειρική σταθερά $\alpha_{cp}$ είναι υπεύθυνη για την προσαρμογή σε πραγματικά δεδομένα και την πιθανή διαφοροποίηση
μεταξύ των εμπρόσθιων και οπίσθιων τροχών.

Το μήκος της περιοχής προσκόλλησης $C_{\text{πρ.},l}$ δίνεται από τον λόγο $C_s(\alpha)$ 
που εκφράζει το ποσοστό του μήκους του πέλματος που είναι προσκολλημένο. Ο συντελεστής $C_s(\alpha)$
λαμβάνεται από γραμμική παρεμβολή μεταξύ δύο γνωστών καταστάσεων (γωνιών ολίσθησης).

\begin{equation}
    C_s(\alpha) = \dfrac{\alpha}{\alpha_c}(C_{s2}-C_{s1}) + C_{s1}
    \label{eq:contactArea2}
\end{equation}

Επομένως, τελικά το εμβαδό της επιφάνειας επαφής του πέλματος για εμπρός και πίσω τροχούς δίνεται από

\begin{equation}
    \begin{aligned}
        A_{\text{πέλμα},f} &= C_{wf}C_sC_l\\
        A_{\text{πέλμα},r} &= C_{wr}C_sC_l
    \end{aligned}
    \label{eq:contactArea3}
\end{equation}

\subsubsection{Εξίσωση κατάστασης}

Η θερμική εξίσωση κατάστασης προκύπτει από ισοζύγιο ενέργειας στη μάζα του πέλματος $m_c$.

\begin{equation}
    \frac{dT}{dt} = \frac{Q_1 + Q_2 - Q_3 - Q_4}{m_c \cdot c_{p,c}}
    \label{eq:dTdt_full}
\end{equation}

Εδώ $Q_1, Q_2$ είναι οι πηγές θερμότητας (\cref{eq:Q1,eq:Q2}) και $Q_3, Q_4$ οι απώλειες
(\cref{eq:Q3,eq:Q4}). Στον παρονομαστή η ποσότητα $m_c \cdot c_{p,c}$ εκφράζει τη θερμική
αδράνεια: μεγαλύτερη μάζα σημαίνει βραδύτερη απόκριση.\\[0.5cm]

Οι προσομοιωτές γύρου λειτουργούν σε πεδίο θέσης $s$ (κατά μήκος πίστας) και όχι χρόνου.
Εφαρμόζοντας τον κανόνα αλυσίδας $dt = ds/v$:

\begin{equation}
    \frac{dT}{ds} = \frac{1}{v}\cdot\frac{dT}{dt} = \frac{Q_1 + Q_2 - Q_3 - Q_4}{m_c \cdot c_{p,c} \cdot v}
    \label{eq:dTds_full}
\end{equation}

Αξίζει να παρατηρηθεί η σημασία της διαίρεσης με $v$: σε χαμηλή ταχύτητα,
ακόμα και μικρό $Q_{\rm net}$ οδηγεί σε μεγάλο $dT/ds$ -- γι' αυτό οι αργές στροφές μπορεί
να είναι πιο επιβαρυντικές για το ελαστικό ανά μέτρο πίστας. Αν αυτό συνδυαστεί με στροφή
όπου μπορεί να έχουμε μεγάλη παραγωγή θερμότητας, αυτό ερμηνεύεται φυσικά ως το όχημα να
ξοδεύει περισσότερο χρόνο σε κατάσταση που παράγεται πολλή θερμότητα.

% ------------------------------------------------------------
%  3) Θερμικό μοντέλο (τιμές τροχού FL)
% ------------------------------------------------------------
\begin{table}[htbp]
    \centering
    \small
    \renewcommand{\arraystretch}{1.2}
    \begin{tabular}{@{}l l l@{}}
        \toprule
        \textbf{Παράμετρος} & \textbf{Τιμή (Μονάδα)} & \textbf{Περιγραφή} \\
        \midrule
        $p_1$                & $0.678$                 & Γραμμικός συντελεστής όρου τριβής $Q_1$ \\
        $p_2$                & $5.12 \cdot 10^{-3}$    & Γραμμ. συντ. $F_x$ όρου έργου παραμόρφωσης $Q_2$ \\
        $p_3$                & $1.16 \cdot 10^{-2}$    & Γραμμ. συντ. $F_y$ όρου έργου παραμόρφωσης $Q_2$ \\
        $p_4$                & $2.89 \cdot 10^{-2}$    & Γραμμ. συντ. $F_z$ όρου έργου παραμόρφωσης $Q_2$ \\
        $p_5$                & $1.09 \cdot 10^{-2}$    & Γραμμικός συντελεστής όρου συναγωγής $Q_3$ \\
        $p_6$                & $2.18$                  & Εκθέτης ταχύτητας όρου συναγωγής $Q_3$ \\
        $C_{wf}$               & $0.20$ m                & Πλάτος εμπρόσθιου ελαστικού \\
        $C_{wr}$               & $0.22$ m                & Πλάτος οπίσθιου ελαστικού \\
        $h_t$                & $12$ kW/m$^2$K          & Συντ. αγωγιμότητας πέλματος-οδοστρώματος \\
        $m_c$                & $0.5$ kg                & Ενεργός θερμική μάζα πέλματος \\
        $c_{p,c}$            & $1800$ J/kg K           & Ειδική θερμοχωρητικότητα ελαστομερούς \\
        $\alpha_{cp}$        & $0.056$                & Σταθερά παραμόρφωσης πέλματος με κατακόρυφο φορτίο \\
        $\alpha_c$           & $8^\circ$           & Γωνία ολίσθησης αναφοράς για γραμμ. παρ. \\
        $C_{s1}$             & $0.3$                  & Ποσοστό περιοχής ολίσθησης -- σημείο 1\\
        $C_{s2}$             & $0.8$                  & Ποσοστό περιοχής ολίσθησης -- σημείο 2\\
        $T_{\mathrm{amb}}$   & $25\,^\circ\mathrm{C}$  & Θερμοκρασία αέρα περιβάλλοντος \\
        $T_{\mathrm{track}}$ & $35\,^\circ\mathrm{C}$  & Θερμοκρασία επιφάνειας οδοστρώματος \\
        \bottomrule
    \end{tabular}
    \caption{Παράμετροι θερμικού μοντέλου. Παρατίθενται οι τιμές του εμπρόσθιου
    αριστερού τροχού ως ενδεικτικές.}
    \label{tab:params_thermal}
\end{table}